\chapter{KONSTRUKSI GELFAND-NAIMARK-SEGAL}\label{gelfand_naimark_segal}

\section{Fungsional Linear Positif}
\begin{defn} Misalkan $A$ suatu aljabar dengan involusi. Untuk setiap fungsional linear $\tau$ pada
$A$ dan setiap $a \in A$ definisikan
 \index{<unopurstar@$\tau^\star$ (konjugat dari $\tau$)}%
$\tau^\star(a) = \conj{\tau(a^*)}$.  Kita mengatakan bahwa $\tau$ bersifat
 \index{Hermitian!fungsional linear}%
 \index{linear!fungsional!Hermitian}%
\df{Hermitian} jika $\tau^\star = \tau$.  Perhatikan bahwa suatu fungsional linear $\tau\colon A
\sto \C$ bersifat Hermitian jika dan hanya jika mempertahankan involusi; yaitu, $\tau^\star =
\tau$ jika dan hanya jika $\tau(a^*) = \conj{\tau(a)}$ untuk semua $a \in A$.
\end{defn}

\begin{cau} $\tau^\star$ yang didefinisikan di atas jangan dikacaukan dengan pemetaan adjoin
biasa $\tau^*\colon \C^* \sto A^*$. Konteks (atau penggunaan kaca pembesar) semestinya
memperjelas mana yang dimaksud.
\end{cau}

\begin{prop} Suatu fungsional linear $\tau$ pada aljabar-$C^*$ $A$ bersifat Hermitian jika dan hanya jika
$\tau(a)~\in~\R$ setiap kali $a$ swaadjoin.
\end{prop}

Seperti halnya selalu terjadi pada pemetaan di antara ruang-ruang vektor terurut, pemetaan
\emph{positif} ialah pemetaan yang membawa elemen positif ke elemen positif.

\begin{defn} Suatu fungsional linear $\tau$ pada aljabar-$C^*$ $A$ bersifat
 \index{positif!fungsional linear}%
 \index{linear!fungsional!positif}%
\df{positif} jika $\tau(a) \ge 0$ untuk setiap $a \in A$ yang memenuhi $a \ge \vc 0$.
\end{defn}


\begin{prop} Setiap fungsional linear positif pada suatu aljabar-$C^*$ bersifat Hermitian.
\end{prop}

\begin{prop} Keluarga fungsional linear positif merupakan suatu kerucut konveks proper dalam ruang
vektor real semua fungsional linear Hermitian pada suatu aljabar-$C^*$. Kerucut itu menginduksi suatu
urutan parsial pada ruang vektor tersebut: $\tau_1 \le \tau_2$ setiap kali $\tau_2 - \tau_1$
bersifat positif.
\end{prop}

\begin{defn} Suatu
 \index{keadaan}%
\df{keadaan} pada aljabar-$C^*$ $A$ adalah fungsional linear positif $\tau$ pada $A$ sedemikian sehingga
$\norm\tau = 1$. Jika $A$ beridentitas, syarat ini ekuivalen dengan $\tau(\vc 1_A) = 1$.
\end{defn}

\begin{exam} Misalkan $x$ suatu vektor dalam ruang Hilbert $H$. Definisikan
   \[ \omega_x:\ofml B(H) \sto \C\colon T \mapsto \langle Tx,x \rangle\,. \]
Maka $\omega_x$ merupakan fungsional linear positif pada $\ofml B(H)$. Jika $x$ suatu
vektor satuan, maka $\omega_x$ merupakan keadaan pada~$\ofml B(H)$. Suatu keadaan $\tau$ disebut
 \index{vektor!keadaan}%
 \index{keadaan!vektor}%
\df{keadaan vektor} jika $\tau = \omega_x$ untuk suatu vektor satuan~$x$.
\end{exam}

\begin{prop}[Ketaksamaan Schwarz]\label{0025105} Jika $\tau$ suatu fungsional linear positif pada
aljabar-$C^*$ $A$, maka
   \[ \abs{\tau(b^*a)}^2 \le \tau(a^*a)\tau(b^*b) \]
untuk semua $a$, $b \in A$.
 \index{ketaksamaan Schwarz}%
 \index{ketaksamaan!Schwarz}%
\end{prop}

\begin{prop} Suatu fungsional linear $\tau$ pada aljabar-$C^*$ beridentitas $A$ bersifat positif jika dan hanya jika
fungsional itu terbatas dan $\norm\tau = \tau(\vc 1_A)$.
\end{prop}

\begin{proof} Lihat \cite{KadisonR:1983}, halaman 256--257. \ns
\end{proof}



















\section{Representasi}
\begin{defn}\label{0028} Misalkan $A$ suatu aljabar-$C^*$. Suatu
 \index{representasi}%
\df{representasi} dari $A$ adalah pasangan $(\pi,H)$ dengan $H$ ruang Hilbert dan $\pi
\colon A \sto \ofml B(H)$ suatu homomorfisme-$*\,$. Biasanya cukup dikatakan bahwa $\pi$ merupakan
representasi dari~$A$. Apabila kita ingin menekankan peran ruang Hilbert tertentu itu,
kita mengatakan bahwa $\pi$ \emph{merupakan representasi dari $A$ pada~$H$.} Bergantung pada konteks, kita
dapat menulis $\pi_a$ atau $\pi(a)$ untuk operator ruang Hilbert yang merupakan citra
elemen aljabar $a$ di bawah~$\pi$. Suatu representasi $\pi$ dari $A$ pada $H$ disebut
 \index{representasi!tak terdegenerasi}%
 \index{tak terdegenerasi!representasi}%
\df{tak terdegenerasi} jika $\pi(A)H$ padat di~$H$.
\end{defn}

 \index{konvensi!representasi aljabar-$C*$ beridentitas bersifat beridentitas}%
\begin{conv} Kita menambahkan persyaratan berikut pada definisi sebelumnya: jika aljabar-$C^*$
$A$ beridentitas, maka suatu representasi dari $A$ harus berupa homomorfisme-$*\,$ \emph{beridentitas}.
\end{conv}

\begin{defn} Suatu representasi $\pi$ dari aljabar-$C^*$ $A$ pada ruang Hilbert $H$ disebut
 \index{representasi!setia}%
 \index{representasi setia}%
\df{setia} jika representasi itu injektif. Jika terdapat suatu vektor $x \in H$ sedemikian sehingga
$\pi^{\sto}(A)x = \{\pi_a(x)\colon a \in A\}$ padat di $H$, maka kita mengatakan bahwa
representasi $\pi$ bersifat
 \index{representasi!siklik}%
 \index{siklik!representasi}%
\df{siklik} dan bahwa $x$ merupakan
 \index{representasi!vektor siklik bagi suatu}%
 \index{siklik!vektor}%
\df{vektor siklik} bagi~$\pi$.
\end{defn}

\begin{exam}\label{002802} Misalkan $(S, \sfml A, \mu)$ suatu ruang ukur $\sigma$-hingga dan
$L_\infty = L_\infty(S, \sfml A, \mu)$ aljabar-$C^*$ dari fungsi-fungsi yang terbatas secara esensial dan
terukur terhadap $\mu$ pada~$S$. Seperti yang kita lihat dalam contoh~\ref{C063527}, untuk setiap $\phi \in
L_\infty$ operator perkalian $M_\phi$ yang bersesuaian merupakan operator pada
ruang Hilbert $L_2 = L_2(S, \sfml A, \mu)$. Pemetaan $M\colon L_\infty \sto \ofml
B(L_2)\colon \phi \mapsto M_\phi$ merupakan representasi setia dari aljabar-$C^*$
$L_\infty$ pada ruang Hilbert~$L_2$.
\end{exam}

\begin{exam} Misalkan $\fml C([0,1])$ aljabar-$C^*$ dari fungsi-fungsi kontinu
pada interval~$[0,1]$. Untuk setiap $\phi \in \fml C([0,1])$, operator perkalian
$M_\phi$ yang bersesuaian merupakan operator pada ruang Hilbert $L_2 = L_2([0,1])$
dari fungsi-fungsi pada~$[0,1]$ yang terintegralkan kuadrat terhadap ukuran Lebesgue.
Pemetaan $M\colon \fml C([0,1]) \sto \ofml B(L_2)\colon \phi \mapsto M_\phi$ merupakan
representasi setia dari aljabar-$C^*$ $\fml C([0,1])$ pada ruang Hilbert~$L_2$.
\end{exam}

\begin{exam} Andaikan $\pi$ suatu representasi dari aljabar-$C^*$ beridentitas $A$ pada
ruang Hilbert~$H$ dan $x$ suatu vektor satuan dalam~$H$. Jika $\omega _x$ merupakan
keadaan vektor yang bersesuaian pada $\ofml B(H)$, maka $\omega_x \circ \pi$ merupakan keadaan pada~$A$.
\end{exam}

\begin{exer} Misalkan $X$ suatu ruang Hausdorff kompak lokal. Carilah suatu representasi
isometrik (dan karena itu setia) $(\pi,H)$ dari aljabar-$C^*$ $C_0(X)$ pada suatu
ruang Hilbert~$H$.
\end{exer}

\begin{defn} Misalkan $\rho$ suatu keadaan pada aljabar-$C^*$ $A$. Maka
   \[ L_\rho := \{a \in A\colon \rho(a^*a) = 0\} \]
disebut
 \index{kiri!kernel}%
 \index{kernel!kiri}%
\df{kernel kiri} dari~$\rho$.
\end{defn}

Ingat bahwa sebagai bagian dari pembuktian \emph{ketaksamaan Schwarz}~\ref{0025105} bagi fungsional
linear positif, kita telah memverifikasi hasil berikut.

\begin{prop} Jika $\rho$ suatu keadaan pada aljabar-$C^*$ $A$, maka $\langle a,b \rangle_0 :=
\rho(b^*a)$ mendefinisikan suatu hasil kali dalam semu pada~$A$.
\end{prop}

\begin{cor} Jika $\rho$ suatu keadaan pada aljabar-$C^*$ $A$, maka kernel kirinya $L_\rho$ merupakan suatu
subruang vektor dari $A$ dan $\langle\,[a],[b]\,\rangle := \langle a,b \rangle_0$ mendefinisikan
suatu hasil kali dalam pada ruang vektor hasil bagi~$A/L_\rho$.
\end{cor}

\begin{prop} Misalkan $\rho$ suatu keadaan pada aljabar-$C^*$ $A$ dan $a \in L_\rho$. Maka $\rho(b^*a)
= 0$ untuk setiap $b \in A$.
\end{prop}

\begin{prop} Jika $\rho$ suatu keadaan pada aljabar-$C^*$ $A$, maka kernel kirinya $L_\rho$ merupakan suatu
ideal kiri tertutup dalam~$A$.
\end{prop}


























\section{Konstruksi GNS dan Teorema Gelfand-Naimark Ketiga}
Teorema berikut dikenal sebagai
 \index{Gelfand-Naimark-Segal!konstruksi}%
 \index{konstruksi!Gelfand-Naimark-Segal}%
\emph{konstruksi Gelfand-Naimark-Segal} (\emph{konstruksi GNS}).

\begin{thm}[Konstruksi GNS]\label{0029} Misalkan $\rho$ suatu keadaan pada aljabar-$C^*$ $A$. Maka
terdapat suatu representasi siklik $\pi_\rho$ dari $A$ pada ruang Hilbert $H_\rho$ dan suatu
vektor satuan siklik $x_\rho$ untuk $\pi_\rho$ sedemikian sehingga $\rho = \omega_{x_\rho} \circ
\pi_\rho$.
\end{thm}

\begin{notn} Dalam pembahasan berikut, $\pi_\rho$, $H_\rho$, dan $x_\rho$ adalah representasi
siklik, ruang Hilbert, dan vektor satuan siklik yang keberadaannya dijamin oleh konstruksi
GNS ketika diterapkan pada suatu keadaan $\rho$ yang diberikan pada suatu aljabar-$C^*$ $A$.
\end{notn}

\begin{prop} Misalkan $\rho$ suatu keadaan pada aljabar-$C^*$ $A$ dan $\pi$ suatu representasi siklik
dari $A$ pada ruang Hilbert $H$ sedemikian sehingga $\rho = \omega_x \circ \pi$ untuk suatu vektor
satuan siklik $x$ bagi~$\pi$. Maka terdapat suatu pemetaan uniter $U$ dari $H_\rho$ ke~$H$ sedemikian sehingga
$x = Ux_\rho$ dan $\pi(a) = U \pi_\rho(a) U^*$ untuk semua $a \in A$.
\end{prop}

\begin{defn} Misalkan $\{H_\lambda\colon \lambda \in \Lambda\}$ suatu keluarga ruang Hilbert.
Nyatakan dengan $\bigoplus\limits_{\lambda \in \Lambda} H_\lambda$ himpunan semua fungsi
$x\colon \Lambda \sto \bigcup\limits_{\lambda \in \Lambda} H_\lambda\colon \lambda
\mapsto x_\lambda$ sedemikian sehingga $x_\lambda \in H_\lambda$ untuk setiap $\lambda \in \Lambda$ dan
$\D\sum\limits_{\lambda \in \Lambda} \norm{x_\lambda}^2 < \infty$. Pada
$\bigoplus\limits_{\lambda \in \Lambda} H_\lambda$, definisikan penjumlahan dan perkalian
skalar secara titik demi titik; yaitu, $(x + y)_\lambda = x_\lambda + y_\lambda$ dan $(\alpha
x)_\lambda = \alpha x_\lambda$ untuk semua $\lambda \in \Lambda$, dan definisikan hasil kali dalam
dengan $\langle x,y \rangle = \sum_{\lambda \in \Lambda} \langle x_\lambda,y_\lambda\rangle$.
Operasi-operasi ini terdefinisi dengan baik dan menjadikan $\bigoplus\limits_{\lambda \in \Lambda}
H_\lambda$ suatu ruang Hilbert. Ruang ini adalah
 \index{jumlah langsung!ruang Hilbert}%
\df{jumlah langsung} dari ruang-ruang Hilbert~$H_\lambda$. Berbagai notasi untuk elemen-elemen jumlah
langsung ini digunakan dalam literatur: $x$, $(x_\lambda)$, $(x_\lambda)_{\lambda \in
\Lambda}$, dan $\oplus_\lambda x_\lambda$ lazim digunakan.

Sekarang andaikan bahwa $\{T_\lambda\colon \lambda \in \Lambda\}$ suatu keluarga operator ruang
Hilbert dengan $T_\lambda \in \ofml B(H_\lambda)$ untuk setiap $\lambda \in \Lambda$.
Andaikan pula bahwa $\sup\{\norm{T_\lambda}\colon \lambda \in \Lambda\} < \infty$. Maka
$T(x_\lambda)_{\lambda \in \Lambda} = (T_\lambda x_\lambda)_{\lambda \in \Lambda}$
mendefinisikan suatu operator pada ruang Hilbert $\bigoplus_\lambda H_\lambda$. Operator $T$
biasanya dinotasikan dengan $\bigoplus_\lambda T_\lambda$ dan disebut
 \index{jumlah langsung!operator ruang Hilbert}%
\df{jumlah langsung} dari operator-operator~$T_\lambda$.
\end{defn}

\begin{prop} Pernyataan-pernyataan dalam definisi sebelumnya bahwa
\smash[b]{$\bigoplus\limits_{\lambda \in \Lambda} H_\lambda$} merupakan ruang Hilbert dan
$\bigoplus_\lambda T_\lambda$ merupakan operator pada $\bigoplus\limits_{\lambda \in \Lambda}
H_\lambda$ adalah benar.
\end{prop}

\begin{exam} Misalkan $A$ suatu aljabar-$C^*$ dan $\{\pi_\lambda\colon \lambda \in \Lambda\}$ suatu
keluarga representasi dari $A$ pada ruang-ruang Hilbert $H_\lambda$ sedemikian sehingga $\pi_\lambda(a)
\in \ofml B(H_\lambda)$ untuk setiap $\lambda \in \Lambda$ dan setiap $a \in A$. Maka
  \[ \pi = \textstyle{ \bigoplus_\lambda \pi_\lambda\colon A \sto
       \ofml B\bigl(\bigoplus_\lambda H_\lambda\bigr)
       \colon a \mapsto \bigoplus_\lambda \pi_\lambda(a)} \]
merupakan representasi dari $A$ pada ruang Hilbert $\bigoplus_\lambda H_\lambda$. Representasi ini adalah
 \index{jumlah langsung!representasi}%
\df{jumlah langsung} dari representasi-representasi~$\pi_\lambda$.
\end{exam}

\begin{thm}\label{thm_exist_faith_rep} Setiap aljabar-$C^*$ memiliki suatu representasi setia.
\end{thm}

\begin{proof} Lihat \cite{Blackadar:2006}, Korolari II.6.4.10; \cite{Conway:1990}, halaman 253;
\cite{DoranB:1986}, Teorema 19.1; \cite{Fillmore:1996}, Teorema 5.4.1;
\cite{KadisonR:1983}, halaman 281; atau \cite{Murphy:1990}, Teorema 3.4.1. \ns
\end{proof}

Suatu perumusan ulang yang jelas dari teorema sebelumnya adalah versi ketiga dari
\emph{teorema Gelfand-Naimark}, yang menyatakan bahwa setiap aljabar-$C^*$ (pada hakikatnya)
merupakan suatu aljabar operator ruang Hilbert.

\begin{cor}[Teorema Gelfand-Naimark III]\label{002973} Setiap aljabar-$C^*$
 \index{Gelfand-Naimark!Teorema III}
isomorfik-$*\,$ secara isometrik dengan suatu subaljabar-$C^*$ dari $\ofml B(H)$ untuk suatu ruang
Hilbert~$H$.
\end{cor}


\endinput
